\documentclass[11pt]{article}

\usepackage{amsmath,amssymb,graphicx}
\usepackage[margin=1in]{geometry}
\usepackage{booktabs}

\title{Delayed Nonlocal Kernel Memory Probe on a Kernel-Induced Cochain Operator Architecture}
\author{Tomislav Rupi\'c \\ Independent Researcher}
\date{March 2026}

\begin{document}

\maketitle

\begin{abstract}
Stage 22 tests the strongest weak-memory hypothesis yet applied to the HAOS-IIP pre-geometric architecture. Earlier Phase II probes established that local scalar reinforcement, mesoscopic basin-residence reinforcement, transport-corridor reinforcement, directional flux-driven reinforcement, mixed scalar-vector reinforcement, and mutual registration gating all remain persistence-null despite bounded nonzero activity. Stage 22 therefore asks whether persistence might require not stronger instantaneous interaction, but temporally accumulated and spatially distributed kernel memory. The tested law class keeps the operator architecture fixed while allowing the kernel weights to depend on a delayed nonlocal functional of past field activity. Three memory laws are evaluated using the same canonical representatives as the late pre-geometric atlas: a single-lag retarded sample, an exponential-history kernel, and a finite-window historical average. All three use the same smoothed nonlocal energy-density driver, the same weak coupling ladder $\varepsilon\in\{0.01,0.02\}$, and the same persistence-first promotion criteria. The 27-run base pilot is uniformly negative. No composite lifetime gain appears, no binding persistence gain appears, no coarse basin persistence gain appears, no corridor dwell gain appears, no transport locking appears, and no refinement promotion is earned. The negative remains structurally informative because every non-null memory class is active. The single-lag and finite-window classes reach memory-field norms of approximately $1.35\times10^{-1}$ with kernel deviation norms of approximately $2.7\times10^{-3}$ and local deformation up to $7.8\times10^{-3}$. The exponential-history class produces the strongest response, with memory-field norm up to $2.65\times10^{-1}$, kernel deviation norm up to $5.3\times10^{-3}$, and local deformation up to $1.52\times10^{-2}$. Activation duty cycle reaches $0.6667$ in all active classes. The bounded conclusion is that delayed nonlocal kernel memory, in the tested weak and support-preserving class, is still insufficient to generate anchoring or persistence on this architecture.
\end{abstract}

\section{Why Stage 22 Exists}

Stage 20 and Stage 21 already define a sharp negative plateau. Weak substrate reinforcement across local, mesoscopic, directional, and mixed symmetry classes failed to produce persistence. Selective interaction eligibility also failed, even when spatial and topology-style registration windows became active. The resulting question is therefore narrower than a generic search for ``what might work.'' It is whether temporal accumulation itself is the missing ingredient.

The physical intuition behind Stage 22 is simple. A packet cannot bind if the substrate only reacts where the packet is now, or only to instantaneous compatibility. A delayed nonlocal memory field is the first tested law in the program that can, in principle, construct a coherent historical structure extending beyond the packet's current position. The pilot asks whether that history-weighted field can become the first persistence-relevant anchoring channel.

\section{Tested Memory Class}

Stage 22 keeps the graph connectivity, projector convention, packet family, and persistence ladder fixed. Only the kernel evolution law changes. In executable form the driver is a smoothed nonlocal energy-density signal $\Phi_{ij}(t)$ sampled on the frozen kernel support. The kernel remains support-preserving and symmetric, with a hard clamp on total deformation.

Three delayed-memory laws are tested.
\begin{enumerate}
\item \textbf{Class A: single-lag retarded reinforcement.} The kernel responds to a discrete delayed sample of the driver, $K(t)=K^{(0)}+\varepsilon\Phi(t-\tau)$.
\item \textbf{Class B: exponential-history kernel.} The kernel responds to a finite-horizon exponential accumulation of past driver values.
\item \textbf{Class C: finite-window historical averaging.} The kernel responds to a bounded short-horizon average of the recent driver history.
\end{enumerate}

All three classes are evaluated on the same representative triad:
\begin{itemize}
\item clustered composite anchor,
\item phase-ordered symmetric triad,
\item counter-propagating corridor.
\end{itemize}
Each class is run under null, very weak, and weak coupling, giving a 27-run base pilot at resolution $n=12$.

\section{Persistence-First Result}

The persistence-first read is uniform across the entire Stage 22 matrix. All 27 runs remain in the label ``no delayed memory effect.'' The decisive observables remain exactly flat against class-matched nulls:
\begin{itemize}
\item composite lifetime delta remains $0.0000$,
\item binding persistence delta remains $0.0000$,
\item coarse basin persistence delta remains $0.0000$,
\item corridor dwell delta remains $0.0000$,
\item no new stabilized ordering class appears,
\item no run earns promotion to the $12\rightarrow24$ refinement ladder.
\end{itemize}

This is therefore a regime-level negative. Temporal accumulation does not open a persistence channel in the tested weak bounded class.

\section{Why the Negative Is Real}

The null does not come from a dead memory law. All three classes enter the substrate measurably.

\subsection{Class A: Single-Lag Retarded Reinforcement}

The single-lag class reaches memory-field norms up to $0.1354$, kernel deviation norms up to $2.7\times10^{-3}$, and local deformation up to $7.7\times10^{-3}$. Activation duty cycle reaches $0.6667$ in the active runs. The clustered composite anchor, the symmetric triad, and the corridor control all see measurable delayed field structure. None convert that into persistence.

\subsection{Class B: Exponential-History Kernel}

The exponential-history class is the strongest in the Stage 22 family. It reaches memory-field norms up to $0.2647$, kernel deviation norms up to $5.3\times10^{-3}$, and local deformation up to $1.52\times10^{-2}$. Activation duty cycle again reaches $0.6667$. This class therefore provides the strongest test of temporal accumulation versus instantaneous response. It still fails completely at the persistence ladder.

\subsection{Class C: Finite-Window Historical Averaging}

The finite-window class closely tracks the single-lag class in scale. It reaches memory-field norms up to $0.1352$, kernel deviation norms up to $2.7\times10^{-3}$, and local deformation up to $7.8\times10^{-3}$. The history field is active and bounded. Again, no persistence consequence follows.

\section{Comparative Hierarchy}

The three delayed-memory laws differ in response strength, but not in outcome.

\begin{table}[h]
\centering
\begin{tabular}{@{}llll@{}}
\toprule
Memory class & Memory-field scale & Kernel deformation scale & Persistence outcome \\
\midrule
Single-lag retarded & $\sim 1.35\times10^{-1}$ & $\sim 2.7\times10^{-3}$ & negative \\
Exponential history & $\sim 2.65\times10^{-1}$ & $\sim 5.3\times10^{-3}$ & negative \\
Finite-window average & $\sim 1.35\times10^{-1}$ & $\sim 2.7\times10^{-3}$ & negative \\
\bottomrule
\end{tabular}
\caption{Stage 22 delayed-memory hierarchy. All three memory laws remain bounded and active, yet none produce composite lifetime gain, binding persistence gain, basin-persistence gain, or corridor locking.}
\end{table}

This hierarchy matters because it rules out an easy under-activation explanation. The exponential-history class is clearly stronger than the single-lag and finite-window classes, yet still leaves every persistence observable unchanged.

\section{Structural Interpretation}

Stage 22 extends the Phase II boundary in an important way. The architecture now resists persistence under:
\begin{itemize}
\item instantaneous weak substrate reinforcement,
\item weak selective interaction gating,
\item and delayed historical reinforcement.
\end{itemize}

The structural reading is therefore stronger than ``memory does not help.'' What the experiment shows is that temporally accumulated substrate memory can be present without becoming an anchor. The memory field is advected through the same transport medium that already failed to produce binding under local and mesoscopic laws. History exists in the substrate, but it still does not close trajectories into a persistent composite regime.

\section{Interpretation Boundary}

Stage 22 does not support emergent geometry, metric structure, or interaction-law closure. It supports only the bounded statement that delayed nonlocal kernel memory, in the tested weak and support-preserving class, does not create persistence. The result therefore applies to a clear region of model space:
\begin{itemize}
\item temporally extended kernel memory,
\item spatially distributed driver field,
\item fixed graph and packet family,
\item no coefficient escalation,
\item persistence-first evaluation unchanged.
\end{itemize}

This still leaves open stronger or qualitatively different mechanisms, but it closes another physically natural weak-law family cleanly.

\section{Conclusion}

Stage 22 is the first direct test of temporal accumulation as a structural degree of freedom in the HAOS-IIP pre-geometric architecture. The question was whether a packet could build a historical potential structure that later packets might encounter as an anchor. The answer, in the tested class, is no.

All three delayed-memory laws remain bounded and active. The exponential-history kernel produces the strongest memory field and the strongest kernel deformation. None of the three classes move the persistence ladder at all. No composite lifetime gain appears, no binding persistence gain appears, no basin-persistence gain appears, and no transport locking appears.

The resulting statement is therefore sharp:

\begin{quote}
Delayed nonlocal kernel memory is insufficient, at the tested weak couplings, to improve composite lifetime, binding persistence, coarse basin persistence, or corridor locking in the current architecture.
\end{quote}

This leaves the program with an even stronger architectural boundary. Persistence has now resisted instantaneous interaction, selective interaction gating, weak substrate reinforcement, and delayed historical reinforcement. The next coherent move is not another nearby memory variant, but a genuinely new degree of freedom, constraint class, or operator sector.

\end{document}
